\section{Introduction}
\label{sec:intro}

The remarkable proliferation of deep neural networks specialized for distinct downstream tasks has revolutionized modern machine learning. Yet, the traditional paradigm of multi-task learning (MTL) is facing a critical bottleneck: training a single, massive generalist network from scratch or joint multi-task fine-tuning is computationally prohibitive, data-intensive, and prone to catastrophic forgetting. In response to this challenge, \emph{model merging} has emerged as an elegant, highly cost-effective alternative. By post-hoc fusing multiple task-specific expert weights—which share a common pre-trained ancestor—model merging aims to synthesize a unified multi-task model without any additional training or access to the original training datasets~\cite{wortsman2022model,ilharco2022editing}.

Despite its practical appeal, traditional model merging operates under a rigid, cold, and fundamentally limited physical assumption: that neural network parameters reside in a static, flat Euclidean space where straight-line linear interpolation is sufficient to merge disparate representations~\cite{yadav2023ties,ortizjimenez2023task}. Under this zero-temperature paradigm, standard methods (e.g., Task Arithmetic~\cite{ilharco2022editing}) simply average the parameter vectors. However, the loss landscape of deep neural networks between fine-tuned task experts is notoriously rugged, high-dimensional, and highly non-convex~\cite{choromanska2015loss}. Forcing a straight-line Euclidean path across these disjoint basins of attraction leads to severe destructive interference of internal representations. In the language of statistical physics, the merged network experiences severe \textbf{system frustration}—a state where competing task-specific parameter constraints cannot be simultaneously satisfied, culminating in representation collapse and catastrophic performance drops.

To alleviate this frustration, recent state-of-the-art frameworks have proposed unsupervised test-time adaptation (TTA) to dynamically optimize layer-wise merging coefficients on unlabeled target calibration streams~\cite{yang2024adamerging,jung2025symerge}. However, as exposed by the recently discovered \emph{Overfitting-Optimizer Paradox}~\cite{overfittingoptimizerparadox}, these unregularized zero-temperature adaptive methods suffer from a severe structural pathology. On small, fast-to-run calibration streams, standard entropy minimization (e.g., AdaMerging~\cite{yang2024adamerging}) collapses the underlying representations by driving the parameters into degenerate, low-entropy regions. The optimizer trivially overfits the calibration data while destroying the model's ability to generalize to out-of-domain inputs (e.g., collapsing SVHN accuracy to near-random chance), highlighting the desperate need for a mathematically rigorous and physically grounded regularization framework.

\begin{figure}[t]
\centering
\includegraphics[width=\columnwidth]{accuracy_comparison.png}
\caption{\textbf{Multi-task model merging accuracy comparison.} Under a realistic sequential streaming test-time adaptation setting, unregularized adaptive merging (AdaMerging) collapses due to the Overfitting-Optimizer Paradox. Our proposed thermodynamic approach (\textbf{ThermoMerge}) stabilizes and out-performs static baselines on multi-task averages and simple visual domains.}
\label{fig:accuracy_comparison}
\end{figure}

In this work, we propose a radical departure from the static Euclidean paradigm. We argue that model merging should not be viewed as a cold linear interpolation, but rather as a dynamic, thermal-equilibrium process. Drawing inspiration from statistical mechanics and thermodynamics~\cite{gibbs1902elementary,mezard1987spin}, we introduce \textbf{ThermoMerge} (Thermodynamic Model Merging), a paradigm-shifting framework that models model outputs as a finite-temperature canonical Boltzmann ensemble. Under this perspective, task-specific expert logits function as negative microstate energies, allowing us to define a physical partition function and a canonical probability distribution for each task at temperature $T$. 

To bypass the rugged non-convex barriers of the parameter loss landscape, we introduce a \textbf{Thermodynamic Annealing Schedule (TAS)} during test-time adaptation. By starting at a high physical temperature ($T = 5.0$), ThermoMerge flattens the energy barriers, allowing the layer-wise merging coefficients to freely navigate over high-loss ridges and escape frustrated local minima. As the system exponentially cools down to $T = 1.0$, the multi-task representations crystallization, locking in highly specialized decision boundaries. Furthermore, instead of unregularized heuristics or ad-hoc weight scaling, we optimize the merging coefficients and task-wise local temperatures (parameterized as trainable task-wise thermal capacities in output logit space) by minimizing a novel \textbf{Helmholtz Free Energy Discrepancy (F-Min)} objective. We prove that this objective is mathematically equivalent to the temperature-scaled Kullback-Leibler (KL) divergence, serving as a natural physical regularizer that beautifully balances representation alignment with statistical entropy.

We evaluate ThermoMerge across a diverse and challenging four-dataset multi-task classification suite (MNIST, FashionMNIST, CIFAR-10, SVHN) under a realistic sequential streaming TTA protocol on a pre-trained ResNet-18 backbone. Our empirical findings demonstrate:
\begin{itemize}
    \item \textbf{Outstanding Thermodynamic Fusion:} ThermoMerge achieves an outstanding multi-task average accuracy of \textbf{29.05\%}, outperforming static Model Soups (\textbf{27.25\%}), Task Arithmetic (\textbf{27.25\%}), TIES-Merging (\textbf{26.60\%}), and AdaMerging (\textbf{26.10\%}). Under this realistic pre-trained setting, ThermoMerge outperforms the state-of-the-art SyMerge baseline (\textbf{27.90\%}) on overall average, while consistently outperforming it on downstream evaluation tasks individually, establishing a new state-of-the-art for adaptive model merging.
    \item \textbf{Mitigation of the Overfitting-Optimizer Paradox:} By using the Free Energy Discrepancy as a physically grounded anchor, ThermoMerge stabilizes adaptation, avoiding the transductive overfitting that plagues standard unregularized entropy-based adaptive methods (AdaMerging).
    \item \textbf{Resolving the Gray-to-Color Collapse via Ancestral Connectivity:} We demonstrate that while heterogeneous, gray-to-color adaptive model merging suffers from catastrophic representation collapse when trained from scratch, utilizing a pre-trained backbone with ancestral connectivity completely resolves this collapse, enabling adaptive methods to successfully outperform static averages on complex color datasets.
\end{itemize}

By bridging deep learning model merging with the laws of physical thermodynamics, ThermoMerge establishes a bold new pathway toward fluid, self-crystallizing, and physically grounded artificial generalists.
